% Generated by eval/scripts/write_magnitude_report.py; do not edit numerical values.

\section{Post-review Analysis of StrongREJECT Score Magnitude}

\label{sec:score-magnitude}

This supplementary analysis was added after inspecting the binary results. It retains the pre-specified analysis and distinguishes the frequency of positive scores from their magnitude. All scored responses, including zeros, contribute to the mean official score. The identity $\overline{s}=\mathrm{ASR}\,\overline{s}_{s>0}$ is descriptive; the positive-only subset differs across models and does not identify a causal severity effect.

\begin{table}[htbp]\centering\small

\caption{Non-aligned-language StrongREJECT frequency and magnitude}\label{tab:score-magnitude}

\begin{tabular}{lrrr}\hline Model & ASR (\%) & Mean score & Positive mean \\\hline

Sagui-7B & 60.39 & 0.377 & 0.625 \\

Sabiá-3 & 8.23 & 0.049 & 0.600 \\

LLaMAntino-2 & 49.62 & 0.259 & 0.521 \\

ANITA-8B & 10.60 & 0.070 & 0.664 \\

GPT-SW3 & 24.04 & 0.132 & 0.550 \\

AI Sweden Llama-3 & 33.40 & 0.194 & 0.581 \\

BgGPT-7B & 73.26 & 0.511 & 0.697 \\

BgGPT-Gemma-2 & 10.77 & 0.077 & 0.712 \\

Llama-3.1 reference & 12.94 & 0.072 & 0.559 \\

\hline\end{tabular}

\par\smallskip\footnotesize The reference is descriptive only; paired inference uses eight targets. Scores lie in $[0,1]$ and are not probabilities of real-world harm.\end{table}

The paired non-aligned panel contains 29598 scored responses, with mean score 0.209 and positive-only mean 0.618. BgGPT-Gemma-2 illustrates relatively infrequent positive scores with substantial conditional magnitude, whereas LLaMAntino-2 produces positive scores more frequently but with lower average magnitude. This comparison concerns observed subsets, not matched successful responses.

\begin{table}[htbp]\centering\small

\caption{Matched pooled GEE coefficients with prompt-clustered 95\% confidence intervals}\label{tab:magnitude-gee}

\begin{tabular}{lrr}\hline Outcome & Distance & SPEC \\\hline

Binary & -0.091 [-0.130, -0.051] & 1.029 [0.916, 1.142] \\

Graded score & -0.112 [-0.142, -0.081] & 0.954 [0.868, 1.040] \\

Positive-only score & -0.114 [-0.140, -0.087] & 0.387 [0.321, 0.453] \\

\hline\end{tabular}\end{table}

The full-panel graded association remains negative for distance and positive for SPEC, as do the separate weak/strong fits. Thus this extension does not reverse the directional conclusions. However, the distance coefficient changes sign when restricting to cells with at least 100 comparable CONS answers: 0.094 [0.057, 0.130]. This restriction also removes the weak Italian model; the reversal is not attributable solely to better measurement.

The IF/CONS decomposition adds nuance. For the graded outcome, standardized IF and CONS coefficients are 1.054 [0.923, 1.185] and 0.453 [0.406, 0.499], respectively; the matched binary coefficients are 0.508 [0.352, 0.664] and 0.571 [0.507, 0.635]. The decomposition has substantial collinearity (IF VIF 38.41; CONS VIF 5.04), limiting claims about separate contributions. Both retain positive coefficients in single-predictor checks, which omit the other substantive predictors. The larger IF coefficient for the graded response suggests that reading-comprehension performance is relevant to the magnitude dimension, but comparison across outcomes is not a formal test of coefficient differences and does not establish a mechanism.

All 114 fitted models, including triggered single-predictor sensitivities, converged with finite estimates and positive-definite robust covariance. The pooled design VIF is 2.46 for distance and 13.51 for SPEC. Full tables, bootstrap intervals, tokenizer-matched fits, original-cohort, questionable-item and translation sensitivities are exported in the versioned magnitude supplement.

These exploratory intervals are conditional on the observed model/language inventory and are not adjusted for multiplicity. Prompt clustering does not capture arbitrary residual dependence across different prompts in the same model--language cell. The positive-only models condition on the outcome, and the existing translation/scoring audit limitations remain. The extension enriches the behavioral description without replacing the crossed-random-intercept primary analysis.
